% ============================================================================
% WORKBOOK (MODIFIED) — Section 6.6: Angle Relationships
% CCSS 7.G.B.5 (VA: SOL extended angle emphasis)
% Practice-focused reinforcement with multi-step algebraic angle problems
% ============================================================================
\section{Angle Relationships}
\topicTitle{Angle Relationships}

\begin{quickReview}{Angle Pairs and Algebra}
\begin{itemize}[leftmargin=*]
    \item \textbf{Complementary angles} add to $90°$.
    \item \textbf{Supplementary angles} add to $180°$.
    \item \textbf{Vertical angles} are formed by two intersecting lines and are always \textbf{equal}.
    \item \textbf{Adjacent angles} share a side and a vertex.
    \item Angles on a straight line (linear pair) add to $180°$.
\end{itemize}

\medskip
\textbf{Example:} Two vertical angles are $(3x + 10)°$ and $(5x - 20)°$. Since vertical angles are equal: $3x + 10 = 5x - 20 \Rightarrow 30 = 2x \Rightarrow x = 15$. Each angle $= 55°$.
\end{quickReview}

% ── Warm-Up ──────────────────────────────────────────────────────────────────
\begin{practiceBox}[funGreen]{\faSun~Warm-Up}

\practiceHeader[funGreen]{Find the missing angle.}

\begin{multicols}{2}
\prob Two angles are complementary. One is $28°$. Find the other. \answerBlank[2cm]
\answerExplain{$62°$}{$90 - 28 = 62°$.}

\prob Two angles are supplementary. One is $115°$. Find the other. \answerBlank[2cm]
\answerExplain{$65°$}{$180 - 115 = 65°$.}

\prob Two vertical angles: one is $73°$. Find the other. \answerBlank[2cm]
\answerExplain{$73°$}{Vertical angles are equal.}

\prob Angles on a line: $40°$ and $x°$. Find $x$. \answerBlank[2cm]
\answerExplain{$140°$}{$180 - 40 = 140°$.}

\prob Complementary angles: one is $45°$. Find the other. \answerBlank[2cm]
\answerExplain{$45°$}{$90 - 45 = 45°$. The angles are equal.}

\prob Supplementary angles: one is $90°$. Find the other. \answerBlank[2cm]
\answerExplain{$90°$}{$180 - 90 = 90°$. Both are right angles.}
\end{multicols}
\end{practiceBox}

% ── Practice A: Algebraic Angle Equations ────────────────────────────────────
\begin{practiceBox}{Algebraic Angle Equations}

\practiceHeader{Set up an equation using the angle relationship and solve.}

\prob Two supplementary angles are $(3x + 15)°$ and $(2x + 5)°$. Find both angles.
\answerExplain{$111°$ and $69°$}{$3x + 15 + 2x + 5 = 180$. $5x + 20 = 180$. $5x = 160$, $x = 32$. Angles: $111°$ and $69°$. Check: $111 + 69 = 180$ \checkmark.}

\prob Two vertical angles are $(4x + 5)°$ and $(6x - 15)°$. Find $x$ and the angle measure.
\answerExplain{$x = 10$; angles are $45°$}{$4x + 5 = 6x - 15$. $20 = 2x$. $x = 10$. Angle: $4(10) + 5 = 45°$.}

\prob Two complementary angles are $(2x - 8)°$ and $(3x + 3)°$. Find both angles.
\answerExplain{$30°$ and $60°$}{$2x - 8 + 3x + 3 = 90$. $5x - 5 = 90$. $5x = 95$, $x = 19$. Angles: $2(19) - 8 = 30°$ and $3(19) + 3 = 60°$.}

\prob A linear pair: $(7x - 4)°$ and $(3x + 24)°$. Find $x$ and both angles.
\answerExplain{$x = 16$; $108°$ and $72°$}{$7x - 4 + 3x + 24 = 180$. $10x + 20 = 180$. $10x = 160$, $x = 16$. Angles: $108°$ and $72°$. Check: $108 + 72 = 180$ \checkmark.}

\end{practiceBox}

% ── Practice B: Multi-Step Angle Problems ────────────────────────────────────
\begin{practiceBox}[funTeal]{Multi-Step Angle Problems}

\practiceHeader[funTeal]{Use multiple angle relationships to solve.}

\prob Two lines intersect. One angle is $(4x - 10)°$. The angle adjacent to it is $(2x + 30)°$. Find all four angles.
\answerExplain{$\approx 96.7°$, $83.3°$, $96.7°$, $83.3°$}{Adjacent angles are supplementary: $4x - 10 + 2x + 30 = 180$. $6x + 20 = 180$. $6x = 160$, $x = \frac{80}{3} \approx 26.67$. Angles: $4(\frac{80}{3}) - 10 = \frac{290}{3} \approx 96.7°$ and $2(\frac{80}{3}) + 30 = \frac{250}{3} \approx 83.3°$. The vertical angles match.}

\prob Three angles meet at a point on a line: $(x + 20)°$, $(2x)°$, and $40°$. Find $x$ and all three angles.
\answerExplain{$x = 40$; angles are $60°$, $80°$, $40°$}{$(x + 20) + 2x + 40 = 180$. $3x + 60 = 180$. $3x = 120$, $x = 40$. Angles: $60°$, $80°$, $40°$. Check: $60 + 80 + 40 = 180$ \checkmark.}

\prob An angle is $20°$ more than $3$ times its complement. Find both angles.
\answerExplain{$72.5°$ and $17.5°$}{Let the complement be $x$. The angle is $3x + 20$. $x + 3x + 20 = 90$. $4x = 70$, $x = 17.5°$. The angle: $3(17.5) + 20 = 72.5°$. Check: $17.5 + 72.5 = 90$ \checkmark.}

\prob An angle is $30°$ less than twice its supplement. Find both angles.
\answerExplain{$110°$ and $70°$}{Let the supplement be $x$. The angle is $2x - 30$. $x + 2x - 30 = 180$. $3x = 210$, $x = 70°$. The angle: $2(70) - 30 = 110°$. Check: $70 + 110 = 180$ \checkmark.}

\end{practiceBox}

% ── Practice C: True or False ────────────────────────────────────────────────
\begin{practiceBox}[funPurple]{True or False?}

\trueOrFalse{Vertical angles are always supplementary.}
\answerExplain{False}{Vertical angles are always \emph{equal}, not supplementary (unless each is $90°$). An angle and its \emph{adjacent} angle are supplementary.}

\trueOrFalse{If two angles are both complementary and equal, each is $45°$.}
\answerExplain{True}{$x + x = 90 \Rightarrow 2x = 90 \Rightarrow x = 45°$.}

\trueOrFalse{Three angles on a straight line must add to $180°$.}
\answerExplain{True}{Angles on a straight line (same side) always sum to $180°$.}

\trueOrFalse{If two lines intersect, all four angles could be different.}
\answerExplain{False}{Vertical angles are equal, so the four angles come in two matching pairs.}

\end{practiceBox}

% ── Word Problems ────────────────────────────────────────────────────────────
\begin{practiceBox}[funBlue]{\faGlobe~Word Problems}

\wordProblem{Two streets cross to form four angles. One angle is $65°$. What are the other three angles?}{degrees}
\answerExplain{$115°$, $65°$, $115°$}{The opposite (vertical) angle is also $65°$. The two adjacent angles are $180 - 65 = 115°$ each.}

\wordProblem{A ramp makes an angle of $(2x + 12)°$ with the ground. The angle between the ramp and the wall is $(3x - 2)°$. These two angles are complementary. Find both angles.}{degrees}
\answerExplain{$44°$ and $46°$}{$2x + 12 + 3x - 2 = 90$. $5x + 10 = 90$. $5x = 80$, $x = 16$. Angles: $2(16) + 12 = 44°$ and $3(16) - 2 = 46°$.}

\wordProblem{A clock's hour and minute hands form a $120°$ angle. What is the supplement of this angle? What is the complement of $30°$?}{degrees}
\answerExplain{Supplement: $60°$; Complement of $30°$: $60°$}{Supplement: $180 - 120 = 60°$. Complement of $30°$: $90 - 30 = 60°$.}

\end{practiceBox}

% ── Challenge ────────────────────────────────────────────────────────────────
\begin{challengeBox}
\prob Two lines intersect. One of the four angles is $(5x + 10)°$. The adjacent angle is $(3x + 2)°$. Find $x$ and all four angles. Then verify that vertical angles are equal and adjacent angles are supplementary. \answerBlank[5cm]
\answerExplain{$x = 21$; angles: $115°$, $65°$, $115°$, $65°$}{$5x + 10 + 3x + 2 = 180$. $8x + 12 = 180$. $8x = 168$, $x = 21$. Angles: $5(21) + 10 = 115°$ and $3(21) + 2 = 65°$. Verticals: $115° = 115°$ \checkmark, $65° = 65°$ \checkmark. Adjacent: $115 + 65 = 180$ \checkmark.}

\prob The ratio of two supplementary angles is $5 : 4$. Find both angles. \answerBlank[3cm]
\answerExplain{$100°$ and $80°$}{Let the angles be $5k$ and $4k$. $5k + 4k = 180$. $9k = 180$, $k = 20$. Angles: $100°$ and $80°$.}
\end{challengeBox}

\encouragement{You've conquered angle relationships --- from simple pairs to multi-step algebra!}
